% ===================================================================
%  TABLICA Nr 9 — Góry. — Uzdrowisko. — Las. — Polowanie.
%  Meme regime que les precedents.
%
%  ICI LA COLONNE SUIT L'IDO ET NON LE FRANCAIS. Le fac-simile
%  francais de ce tableau a perdu la moitie de ses intertitres et en
%  place deux au mauvais endroit (« Deuxieme scene » a la place de
%  « Paysage de montagne », « La Foret. --- La Chasse. » a la place de
%  « L'Excursionniste »). L'ido, lui, a ses douze intertitres a leur
%  rang, et le neerlandais comme le tcheque les suivent. Cette
%  colonne aussi : le tableau qu'on traduit est l'ido.
%
%  MEME CHOSE POUR DEUX ALINEAS QUE LE FRANCAIS ABREGE : au 5 de la
%  deuxieme scene il oublie la fuite du braconnier et le gibier reste
%  sur l'herbe (lievre, biche, chevreuil, sanglier — renvois 18 a
%  21), et au 8 il refait la phrase du charbonnier. On traduit l'ido,
%  qui les a.
%
%  ET TROIS RENVOIS DE CE TABLEAU SONT CORRIGES SUR LA PAGE DE
%  LECTURE, DANS LA COLONNE FRANCAISE SEULE : elle imprime (11), (14)
%  et (19) la ou l'ido lit (41), (44) et (49). La correction vit dans
%  gravuri/korekti.json et ne touche aucune source. Cette colonne
%  ecrit les numeros de l'ido.
%
%  « pastoro » n'etait pas le pasteur au tableau 7 ; ici « defileo »
%  n'est pas un defile mais un col — przelecz —, et « silvo » (25) la
%  futaie, le vieux peuplement, non la foret en general, qui est
%  « foresto ». La planche montre les deux.
% ===================================================================

%%K t09-apar-1 apar
\VUsaut{4.0mm}
\VUtitre{15.0pt}{60}{TABLICA Nr 9}
\VUsaut{3.0mm}

%%K t09-tit-1 sub
\VUcentre{12.0pt}{0}{\textit{Scena pierwsza.}}
\VUsaut{3.0mm}

%%K t09-tit-2 sub
\VUpk{11.4pt}{40}{Góry. --- Uzdrowisko.}
\VUsaut{3.0mm}

%%K t09-c1-01-1 p
1. --- Od ośmiu dni jestem w Pratz. Nie potrafię wyrazić całego podziwu, jaki poczułem,
kiedy stanąłem przed cudownym \VUgras{łańcuchem górskim} \textsuperscript{(1)} Pirenejów,
którego \VUgras{grzbiet} \textsuperscript{(2)} rysuje się na błękitnym niebie jak olbrzymi
feston.

%%K t09-c1-02-1 p
2. --- Nie wyobrażałem sobie, że pirenejskie \VUgras{szczyty} \textsuperscript{(3)}
osiągają tak wielką \VUgras{wysokość}, i stałem oszołomiony wobec przepięknych
\VUgras{lodowców} \textsuperscript{(5)} i ośnieżonych \VUgras{turni} \textsuperscript{(4)},
po których staczają się tak straszliwe \VUgras{lawiny}.

%%K t09-tit-3 sub
\VUsaut{3.0mm}
\VUpk{10.2pt}{40}{Krajobraz górski.}
\VUsaut{3.0mm}

%%K t09-c1-03-1 p
3. --- Już nazajutrz po przyjeździe poszedłem do starego wygasłego \VUgras{wulkanu}
\textsuperscript{(6)}, który jest blisko Pratz. Na jałowych i ogołoconych brzegach
\VUgras{krateru} widać jeszcze dobrze ślady przejścia \VUgras{lawy}, która płynęła przed
wieloma wiekami po \VUgras{zboczu góry} \textsuperscript{(7)}. Udało mi się znaleźć na
jednym ze \VUgras{stoków} kilka odłamków tej lawy.

%%K t09-c1-04-1 p
4. --- Pratz leży na granicy, a \VUgras{twierdza} \textsuperscript{(8)}, która broni
\VUgras{przełęczy} \textsuperscript{(27)} oddzielającej Francję od Hiszpanii, całkiem
przypomina te stare zamki znad Renu, które zdają się wypatrywać \VUgras{łupu} ze szczytu
swojej skały.

%%K t09-c1-05-1 p
5. --- Ogromny \VUgras{orzeł} \textsuperscript{(9)}, który ma swoje gniazdo niedaleko
\VUgras{fortu} \textsuperscript{(8)}, często przyciąga moją uwagę. Ten \VUgras{ptak
drapieżny}, którego \VUgras{dziób} jest silny, często przynosi swoim małym jagnię albo
koźlę, które trzyma w swoich \VUgras{szponach}. Tak wielkie są jego \VUgras{skrzydła}, że
potrafi długo szybować ze swoim ciężkim brzemieniem.

%%K t09-tit-4 sub
\VUsaut{3.0mm}
\VUpk{10.2pt}{40}{Wycieczkowicz.}
\VUsaut{3.0mm}

%%K t09-c1-06-1 p
6. --- Najprzyjemniejszą przechadzką jest tam wycieczka do wodospadu
\textsuperscript{(10)} w «Kau». \VUgras{Wodospad} spada wirując, a potem znika jako piękny
\VUgras{strumyk} w \VUgras{lesie jodłowym} \textsuperscript{(11)} położonym na zachód od
miasta. Wspięliśmy się przez ten las aż do \VUgras{obserwatorium meteorologicznego}
\textsuperscript{(12)}, a ze szczytu \VUgras{punktu widokowego} dostrzegliśmy całą
\VUgras{panoramę} okolicy. \VUgras{Krajobraz} jest cudowny, a \VUgras{przyroda} zdaje się
szczodrze obdarzać ten kraj wszystkimi skarbami, którymi rozporządza.

%%K t09-c1-07-1 p
7. --- Aby zejść, skorzystaliśmy z \VUgras{kolejki linowej} \textsuperscript{(13)} i
zatrzymaliśmy się na małej \VUgras{stacji} blisko \VUgras{ruin} \textsuperscript{(14)}
starego \VUgras{zamku feudalnego}, zamieszkiwanego niegdyś przez panów na Pratz.

%%K t09-c1-08-1 p
8. --- Biedny zamek! Co też pewnie myśli, widząc w dole zalotne \VUgras{uzdrowisko}
\textsuperscript{(15)}, które jest teraz modnym \VUgras{zdrojem termalnym}?

%%K t09-c1-08-2 p
\VUgras{Klimat} jest tak przyjemny, \VUgras{woda mineralna} tak skuteczna, że
\VUgras{kuracja} stosowana w \VUgras{sanatorium} \textsuperscript{(17)} jest prawdziwą
przyjemnością.

%%K t09-tit-5 sub
\VUsaut{3.0mm}
\VUpk{10.2pt}{40}{Przyjemne miasto zdrojowe.}
\VUsaut{3.0mm}

%%K t09-c1-09-1 p
9. --- Zresztą zrobiono wszystko, żeby uprzyjemnić pobyt. Wielki \VUgras{wiadukt}
\textsuperscript{(16)} zbudowano, aby umożliwić kolej żelazną; \VUgras{park}
\textsuperscript{(18)} \VUgras{zakładu zdrojowego}, w którym się \VUgras{koncertuje},
urządzono w uroczy sposób. Kilka wdzięcznych «\VUgras{willi}» \textsuperscript{(19)}
zbudowano wokół kasyna \textsuperscript{(21)}, a bardzo dobrze utrzymana \VUgras{szosa}
\textsuperscript{(22)} znacznie ułatwia komunikację.

%%K t09-c1-10-1 p
10. --- Zwykła \VUgras{skarpa} \textsuperscript{(23)} oddziela \VUgras{drogę}
\textsuperscript{(20)} od właściwej \VUgras{doliny} \textsuperscript{(25)}, na której dnie
leży wdzięczne małe \VUgras{jezioro} \textsuperscript{(26)}, zasilające przejrzysty
\VUgras{strumień} \textsuperscript{(24)}.

%%K t09-c1-11-1 p
11. --- Po drugiej stronie doliny eksploatuje się \VUgras{kopalnię żelaza}
\textsuperscript{(28)}. Wiele razy chodziłem patrzeć, jak \VUgras{górnicy} zjeżdżają do
\VUgras{szybu}, a nawet wszedłem do \VUgras{chodnika}, w którym spędzają swój dzień,
wydobywając \VUgras{rudę}, która zawiera \VUgras{metal}.

%%K t09-c1-12-1 p
12. --- Ważną \VUgras{hutę} zbudowano blisko kopalni, żeby od razu wykorzystać bogactwa,
które się z niej wydobywa. \VUgras{Wielki piec} \textsuperscript{(29)}, ogrzewany
\VUgras{węglem kamiennym}, którego tu pod dostatkiem, pozwala szybko i tanio otrzymywać
\VUgras{surówkę}.

%%K t09-c1-13-1 p
13. --- Niedaleko kopalni drąży się \VUgras{kamieniołom} \textsuperscript{(30)}. Ani
\VUgras{granitu}, ani \VUgras{porfiru}, ani nawet \VUgras{piaskowca} nie odłupuje stary
\VUgras{kamieniarz} swoim \VUgras{kilofem}, lecz zwykły \VUgras{kamień ciosowy} do budowy
domów.

%%K t09-c1-14-1 p
14. --- Jedną z najpiękniejszych ozdób tego kraju są \VUgras{jodły}
\textsuperscript{(31)}. Wszędzie w głębi \VUgras{wąwozu} \textsuperscript{(32)}, na brzegu
\VUgras{potoku} \textsuperscript{(33)}, \VUgras{przepaści} \textsuperscript{(34)} albo
urwiska ich cienkie igły kładą zielony odcień.

%%K t09-tit-6 sub
\VUsaut{3.0mm}
\VUpk{10.2pt}{40}{Chodzenie pieszo.}
\VUsaut{3.0mm}

%%K t09-c1-15-1 p
15. --- Korzystam z wakacji, żeby \VUgras{długo} spacerować. \VUgras{Drogi}
\textsuperscript{(35)}, które wiją się po górze, pozwalają przebyć, nie zważając na
odległość, kilka \VUgras{kilometrów}, a często kilka \VUgras{mil}.

%%K t09-c1-15-2 p
\VUgras{Słupy graniczne} \textsuperscript{(36)} i \VUgras{drogowskazy}
\textsuperscript{(37)} znaczą drogi, na których często spotyka się młodego
\VUgras{górala} \textsuperscript{(38)} z \VUgras{sakwą} \textsuperscript{(40)} u boku,
niosącego \VUgras{świstaka} \textsuperscript{(39)}, albo jakiegoś \VUgras{Cygana}
\textsuperscript{(41)}, którego \VUgras{futrzana czapa} \textsuperscript{(42)} dziwnie
kontrastuje ze starym podartym \VUgras{kapturem} \textsuperscript{(43)}. Prowadzi on na
\VUgras{łańcuchu} tresowanego \VUgras{niedźwiedzia} \textsuperscript{(44)}.

%%K t09-tit-7 sub
\VUpk{10.2pt}{40}{Wspinaczka górska.}
\VUsaut{3.0mm}

%%K t09-c1-16-1 p
16. --- Prawie co dzień idę z \VUgras{przewodnikiem} \textsuperscript{(48)} ćwiczyć
\VUgras{wspinaczkę} i jestem bardzo dumny, kiedy z \VUgras{plecakiem}
\textsuperscript{(47)} na plecach i «\VUgras{alpensztokiem}» w ręku, jak prawdziwy
\VUgras{turysta} \textsuperscript{(46)}, szturmuję \VUgras{skały} \textsuperscript{(45)}.

%%K t09-c1-17-1 p
17. --- Ze szczytu góry mieszkańcy równiny wyglądają nie grubiej niż mrówki;
\VUgras{stado owiec} \textsuperscript{(50)} pasące się na \VUgras{pastwisku}
\textsuperscript{(49)} zdaje się zabawką dla dzieci. \VUgras{Sosna}
\textsuperscript{(51)} albo i \VUgras{drewniany domek} \textsuperscript{(52)} wygląda
ledwie jak plamka na zieleni.

%%K t09-c1-18-1 p
18. --- Podczas tych wycieczek często spotykamy na \VUgras{ścieżce}
\textsuperscript{(53)} \VUgras{łowcę kozic} \textsuperscript{(54)}, niosącego na ramieniu
zwierzę, które dopiero co zabił.

%%K t09-c1-19-1 p
19. --- Umieściłem w moim zielniku bardzo osobliwą gałązkę \VUgras{paproci}
\textsuperscript{(55)} i piękną różową gałązkę \VUgras{wrzosu} \textsuperscript{(57)},
które zerwałem u wejścia do małej \VUgras{groty} \textsuperscript{(56)} podczas mojego
ostatniego spaceru.

%%K t09-tit-8 sub
\VUsaut{3.0mm}
\VUcentre{12.0pt}{0}{\textit{Scena druga.}}
\VUsaut{3.0mm}

%%K t09-tit-9 sub
\VUpk{11.4pt}{40}{Las. --- Polowanie.}
\VUsaut{3.0mm}

%%K t09-c2-01-1 p
1. --- Wczoraj spędziłem rozkoszny dzień w lesie. Krzątanina, która panuje wśród
zwierząt i \VUgras{owadów} \textsuperscript{(5)} zamieszkujących go, bardzo mnie
zainteresowała.

%%K t09-c2-01-2 p
\VUgras{Pająk} \textsuperscript{(1)}, który tka swoją \VUgras{sieć} \textsuperscript{(2)}
między dwiema gałęziami, żeby złapać przelatującą \VUgras{muchę} \textsuperscript{(3)},
\VUgras{ślimak} \textsuperscript{(4)} mozolnie niosący swoją muszlę, jelonek rogacz
badający swoją zdobycz, pilna mrówka bardzo gorliwa przy \VUgras{mrowisku}
\textsuperscript{(6)}, ci wszyscy malutcy chodzili, przychodzili, pracowali z
nadzwyczajną pilnością.

%%K t09-c2-02-1 p
2. --- \VUgras{Wąż} \textsuperscript{(7)}, którego na pierwszy rzut oka wziąłem za
\VUgras{żmiję}, ale który był tylko zwykłym \VUgras{zaskrońcem}, przyczaił się pod pniem
drzewa i od czasu do czasu wysuwał swoją cienką główkę, która zawsze zdawała się gotowa
ukąsić. Przede mną gruby \VUgras{pomrów} \textsuperscript{(8)} ciężko pełzł, pożarłszy
jedną ze smakowitych \VUgras{poziomek} \textsuperscript{(11)}, które obficie tam rosną.

%%K t09-c2-03-1 p
3. --- Ziemia była usłana \VUgras{kwiatami leśnymi}; \VUgras{pierwiosnki}
\textsuperscript{(9)}, \VUgras{barwinki} \textsuperscript{(10)} i wiele innych roślin,
których nie znałem, kwitły wśród \VUgras{kęp trawy} \textsuperscript{(12)}.

%%K t09-c2-04-1 p
4. --- W tej okolicy \VUgras{trufla} jest prawie tak pospolita jak \VUgras{grzyb}
\textsuperscript{(14)}, a \VUgras{prosię} \textsuperscript{(13)} należące do leśnika
łakomie badało, czy uda mu się znaleźć kilka z nich.

%%K t09-tit-10 sub
\VUsaut{3.0mm}
\VUpk{10.2pt}{40}{Kłusownictwo.}
\VUsaut{3.0mm}

%%K t09-c2-05-1 p
5. --- Byłem świadkiem \VUgras{końca} sceny, która bardzo mnie ubawiła. Dzięki węchowi
swojego \VUgras{psa jamnika} \textsuperscript{(15)} \VUgras{leśnik} \textsuperscript{(16)}
właśnie zaskoczył \VUgras{kłusownika} \textsuperscript{(22)} w chwili, gdy ten szykował
się zabrać \VUgras{zwierzynę} \textsuperscript{(17)}, którą zabił. Woląc porzucić swoją
zdobycz niż dać się aresztować, kłusownik uciekł, a \VUgras{zając} \textsuperscript{(18)},
\VUgras{łania} \textsuperscript{(19)}, \VUgras{sarna} \textsuperscript{(20)} i
\VUgras{dzik} \textsuperscript{(21)} leżały na trawie ku radości leśnika.

%%K t09-c2-06-1 p
6. --- Różnorodność drzew, które zawiera las, jest zadziwiająca. \VUgras{Buki}
\textsuperscript{(23)} i \VUgras{brzozy} \textsuperscript{(26)}, \VUgras{sosny}, które
wydają \VUgras{żywicę} \textsuperscript{(27)}, \VUgras{dęby} \textsuperscript{(29)} i
wiele innych mieszają swoje konary.

%%K t09-c2-06-2 p
\VUgras{Bluszcz} \textsuperscript{(24)}, który przywiera do pnia wysokich drzew, barwi
\VUgras{starodrzew} \textsuperscript{(25)} lśniącą zielenią, która przyjemnie łączy się z
jaśniejszym i blado zielonym odcieniem \VUgras{mchu} \textsuperscript{(31)}, w którym
\VUgras{dzięcioł zielony} \textsuperscript{(30)} szuka owadów.

%%K t09-tit-11 sub
\VUsaut{3.0mm}
\VUpk{10.2pt}{40}{Krajobraz leśny.}
\VUsaut{3.0mm}

%%K t09-c2-07-1 p
7. --- Stare dęby mnie oczarowały. Ich grube skręcone \VUgras{korzenie}
\textsuperscript{(32)}, do połowy wyszłe z ziemi, dają im wspaniały wygląd siły i wigoru,
i pytam siebie, jak \VUgras{właściciel} \textsuperscript{(35)} może godzić się na ich
ścięcie i wydanie \VUgras{pile} \textsuperscript{(34)} \VUgras{drwala}
\textsuperscript{(33)}. Biedne \VUgras{pnie} \textsuperscript{(36)}, obdarte ze swojej
\VUgras{kory} \textsuperscript{(37)} albo rozłupane \VUgras{siekierą}
\textsuperscript{(38)} \VUgras{wyrobnika} \textsuperscript{(39)}, zasmucają mnie.

%%K t09-c2-08-1 p
8. --- Obok stary \VUgras{węglarz} \textsuperscript{(41)} przygotował swoją
\VUgras{łopatą stos} \textsuperscript{(42)} węgla drzewnego, a staruszka przyniosła
\VUgras{wiązkę} \textsuperscript{(40)} \VUgras{martwego drewna} z gałązek ściętych drzew.

%%K t09-tit-12 sub
\VUsaut{3.0mm}
\VUpk{10.2pt}{40}{Polowanie.}
\VUsaut{3.0mm}

%%K t09-c2-09-1 p
9. --- Chciałem pójść odpocząć przez kilka chwil w \VUgras{leśniczówce}
\textsuperscript{(46)}, ale kiedy przybyłem blisko \VUgras{stawu} \textsuperscript{(43)},
na którym pływa piękny \VUgras{łabędź} \textsuperscript{(45)}, zauważyłem, że niedawna
\VUgras{powódź} \textsuperscript{(44)} zmieniła drogę w prawdziwe \VUgras{bagno}.
Musiałem zawrócić i, przechodząc blisko \VUgras{cedru} \textsuperscript{(49)},
\VUgras{topoli} \textsuperscript{(48)} i \VUgras{wiązu} \textsuperscript{(47)}, które są
na skraju lasu, zobaczyłem, jak z \VUgras{zagajnika} \textsuperscript{(54)} wypada
przepiękny \VUgras{jeleń} \textsuperscript{(50)}, ścigany przez zawziętą \VUgras{sforę}
\textsuperscript{(51)}, którą podniecał jeszcze \VUgras{róg} \textsuperscript{(53)}
\VUgras{konnych myśliwych} \textsuperscript{(52)}. Biedne zwierzę było zupełnie
wyczerpane. Padło umierając o kilka kroków ode mnie.

%%K t09-c2-10-1 p
10. --- Syn leśnika, którego przyciągnął zgiełk \VUgras{polowania z psami}, wyszedł z
domu i wróciliśmy we dwóch do lasu. Widział on \VUgras{srokę} \textsuperscript{(56)} na
gałęzi starego dębu. Sądził, że jej gniazdo jest pewnie ukryte w \VUgras{listowiu}
\textsuperscript{(58)}, i chciał to gniazdo złapać.

%%K t09-c2-10-2 p
Zwinny jak \VUgras{wiewiórka} \textsuperscript{(61)}, wspinał się z \VUgras{gałęzi}
\textsuperscript{(55)} na gałąź, ale nie znalazł nic i zdołał tylko spłoszyć starą
\VUgras{wronę} \textsuperscript{(60)} siedzącą na \VUgras{konarze} \textsuperscript{(57)}.
\VUsaut{4.0mm}
\VUfilet{22mm}
